\documentclass[sigconf,natbib=true,anonymous=true]{acmart}

\AtBeginDocument{%
  \providecommand\BibTeX{{\normalfont B\kern-0.5em{\scshape i\kern-0.25em b}\kern-0.8em\TeX}}}

\setcopyright{none}
\settopmatter{printacmref=false}
\renewcommand\footnotetextcopyrightpermission[1]{}
\pagestyle{plain}
\acmConference[FIRE 2026]{Forum for Information Retrieval Evaluation}{December
  2026}{India}

\usepackage{booktabs}
\usepackage{array}
\usepackage{tabularx}
\usepackage{adjustbox}
\usepackage{graphicx}
\usepackage{amsmath}

\begin{document}

\title{Matching and Scoring Checks for Weight-Reusing Encoder Bias Comparisons}

\author{Anonymous Author(s)}

\begin{abstract}
Retrieval systems rank documents, and text classifiers assign labels to documents. Many such systems reuse language encoders, which map text into contextual token representations. When an encoder gives higher probability to text that links a group with a stereotype, later systems may inherit that lexical association. Weight reuse applies the same network block at several depths. A matched training point compares saved model states with similar held-out text loss. This paper asks whether stereotype-association comparisons change when the matched training point and scoring rule change. We train four encoders on the same text. We evaluate three datasets: \emph{CrowS-Pairs}, a paired-sentence stereotype benchmark; \emph{StereoSet}, a sentence-completion stereotype benchmark; and \emph{WinoBias}, a pronoun-resolution capability check. We apply four sensitivity checks. At the lowest-loss matched point on \emph{CrowS-Pairs}, \emph{LoopedBERT} and \emph{HyperloopBERT} have lower per-item stereotype score differences than the unshared \emph{VanillaBERT} baseline after Holm correction, a multiple-comparison adjustment. \emph{HyperloopBERT} also required a lower learning rate, so its result is a bundle of architecture and stability changes. The conclusion does not hold across the checks. In a loss-adjusted \emph{CrowS-Pairs} regression, architecture terms are not distinguishable from zero. Scoring only shared words or using the released scorer, which assigns the stereotypical role from the dataset direction label, does not preserve the corrected result. On \emph{StereoSet}, only the \emph{HyperloopBERT} comparison remains corrected. The \emph{WinoBias} capability check fails, so the results support claims about lexical association only.
\end{abstract}

\ccsdesc{Computing methodologies~Natural language processing}
\ccsdesc{Information systems~Evaluation of retrieval results}
\ccsdesc{Social and professional topics~User characteristics}

\keywords{social bias, masked language models, parameter sharing, evaluation
sensitivity, benchmark reliability, fairness measurement}

\maketitle

\section{Introduction}
\label{sec:intro-bundle}

Information retrieval systems rank documents for user queries, and text
classifiers assign labels to documents. Many such systems start from a
pre-trained language encoder, which maps a text sequence into contextual token
representations. Encoder choice therefore matters before any downstream task is
known. If the encoder prefers text that links a social group with a stereotype,
later systems may inherit that lexical association. Fine-tuning can change the
association, but the initial encoder can still affect model selection.

A masked language model is an encoder trained to predict tokens hidden from
their surrounding words. A social stereotype association is a higher model score
for text that links a social group with a common stereotype. Stereotype
benchmarks estimate such associations by comparing model scores for paired
texts.

One \emph{CrowS-Pairs} item makes the measurement concrete. \emph{CrowS-Pairs}
is a stereotype benchmark of paired sentences that differ in a demographic term.
One item contrasts ``\emph{Cheryl} was bad at saving money and thus had a
financial crisis'' with the same sentence using ``\emph{John}'',
Table~\ref{tab:data}. The task asks which sentence receives the higher model
score. Under the full-sentence rule, each non-boundary token is masked in turn.
The model probability of the original token is recorded. The log probabilities
are averaged. The anti-stereotypical score is subtracted from the stereotypical
score. A positive item score means that the encoder gives the stereotypical
sentence the higher probability.

The practical question is whether one encoder architecture gives lower
stereotype scores than another. An architecture difference cannot be attributed
from one benchmark score, because measured association also changes with model
capability. Model capability here means language-modelling quality, measured by
validation loss on held-out text. Validation loss is the average negative log
probability on held-out text, so lower values indicate better language
modelling. A checkpoint is a saved model state at a training step. A model later
in training can assign more concentrated probabilities on benchmark pairs. A
fixed training budget therefore mixes architecture with training progress.

Prior work gives reasons to be cautious about a single architecture comparison.
Compression studies show that fairness measurements can change when capacity is
changed \citep{hooker2020characterising,xu2022compression,ramesh2023comparative}.
Stereotype-benchmark audits report item-quality problems that weaken conclusions
drawn from one benchmark average
\citep{blodgett2021salmon,wang2025difference}. The remaining gap is whether a
comparison between weight-reusing encoder designs keeps the same sign when
training progress and scoring choices are varied.

This paper runs the \emph{CrowS-Pairs} architecture comparison under matched
points and scoring rules. Weight reuse means that the same block parameters are
applied at more than one depth in the network. Effective depth means the number
of transformer-block applications made to each token. A transformer block is the
repeated attention-and-feed-forward unit used in BERT, the original
bidirectional transformer encoder design \citep{devlin2019bert}. The objective
is to estimate whether lower measured association can be attributed to weight
reuse, rather than to realised training loss or scorer convention.

The comparison varies encoder design while holding the tokeniser, objective,
training text, width, and effective depth fixed. A tokeniser is the fixed
procedure that splits text into model tokens. A lexical association is a model
preference for words or sentences that co-occur with a social group in the
benchmark text. A matched point is a checkpoint selected by crossing the same
validation-loss target. A contrast is the mean paired difference between the
unshared baseline and another encoder. It is computed over the same benchmark
items. Holm correction is a step-down procedure for several contrasts against
one baseline. It controls the family-wise error rate.

The study compares the four encoders in Table~\ref{tab:arch}. \emph{VanillaBERT}
is the unshared baseline with independent transformer blocks. \emph{LoopedBERT}
reuses a small core of blocks across depth. ALBERT is a BERT variant that shares
encoder-block parameters across layers \citep{lan2020albert}.
\emph{ALBERTLoopedBERT} reuses one block across the whole stack, following
ALBERT. \emph{HyperloopBERT} keeps the looped design and adds parallel
representation streams mixed by learned weights \citep{zhu2025hyperconnections}.
A stream is one parallel copy of token representations inside that model.

This paper makes these contributions. The protocol evaluates \emph{CrowS-Pairs},
\emph{StereoSet}, and \emph{WinoBias}. It matches encoders on validation loss
and checks the conclusion across matched points, continuous loss adjustment,
scoring rules, and a second benchmark. The \emph{CrowS-Pairs} result at the
lowest-loss matched point depends on the matched point and the scorer.
\emph{LoopedBERT} shows the same direction at only one matched point,
Figure~\ref{fig:delta}. \emph{LoopedBERT} and \emph{HyperloopBERT} are lower
than \emph{VanillaBERT} only in the single-point corrected comparison,
Table~\ref{tab:band}. The architecture coefficients are not distinguishable
from zero after continuous adjustment for realised loss, Table~\ref{tab:model}.
The alternative scoring rules do not preserve the corrected result,
Table~\ref{tab:rules2}. The category analysis and the failed capability gate
limit the interpretation to lexical association in the masked language modelling
head. The masked language modelling head is the output layer that assigns
probabilities to tokens.

These claims rely on prior work on bias measurement, weight sharing, and
capacity control. The next section states that background and identifies where
the present comparison differs.

\section{Related Work}
\label{sec:related}

Prior work provides the stereotype measurements, the weight-reusing
architectures, and the reason to control capability when comparing encoders.

\emph{CrowS-Pairs} and \emph{StereoSet} compare model scores for stereotypical
and anti-stereotypical text \citep{nangia2020crows,nadeem2021stereoset}. A
pseudo-log-likelihood is the average log probability assigned to original tokens
when each token is masked in turn. Both benchmarks build on masked-position
probing of social associations \citep{kurita2019measuring}.

Benchmark audits report item-quality problems that weaken conclusions drawn
from one average. \citet{blodgett2021salmon} found unclear or inconsistent
items in both benchmarks. \citet{wang2025difference} distinguish
correlation-type benchmarks, which measure association, from tests that measure
whether model behaviour treats groups differently when equal treatment is
expected. The benchmarks used here are correlation-type instruments. The present
study uses the same instruments to ask whether an architecture comparison
changes under alternative scoring choices.

A separate literature changes encoder architecture by reusing weights across
depth. ALBERT ties one encoder block across all layers to reduce storage
\citep{lan2020albert}. \citet{saunshi2025looped} show that a block repeated in a
loop can approach the reasoning performance of a deeper unshared model.
\citet{bae2025mor} learn a per-token recursion depth over a shared block.
\citet{geiping2025huginn}, \citet{zhu2025ouro}, and \citet{frey2026adaptive}
study related recurrent-depth models. A residual connection adds a block input
to the block output before the next layer. Hyper-connections generalise that
connection to several parallel streams with learned mixing
\citep{zhu2025hyperconnections}, and \citet{xie2025mhc} constrain that mixing.
\citet{zeitoun2026hyperloop} combine looping with hyper-connections. The looped
and hyper-connection studies cited above mainly report language quality,
throughput, parameter efficiency, and reasoning accuracy. They do not test
social-bias measurements for these designs.

Capacity and compression studies are the closest precedent. Pruning means
removing model weights after training. \citet{hooker2020characterising} found
that pruning can harm groups unevenly. \citet{xu2022compression} and
\citet{ramesh2023comparative} compare compression methods and fairness
measurements. \citet{voria2026tracing} localise stereotype-carrying neurons
inside pre-trained transformers. The present study varies weight sharing before
training and compares models at matched validation loss.

Fairness in information access is an established concern for retrieval and
ranking systems \citep{ekstrand2022fairness,singh2018fairness}. A claim that
one encoder design carries less stereotype association than another could affect
model selection. The present study contributes measurement evidence behind such
a claim, rather than a new ranking algorithm.

\section{Data}
\label{sec:data}

The data have separate training and evaluation roles. One corpus trains the
encoders, and the evaluation sets measure association or capability. The
concrete examples in Table~\ref{tab:data} make the tasks visible, and the exact
data choices are listed in Table~\ref{tab:datastats}.

All encoders in Table~\ref{tab:arch} are pre-trained on \emph{FineWeb-Edu}, a
filtered subset of web text selected for educational content
\citep{penedo2024fineweb}. The shared corpus size and sequence length are given
in Table~\ref{tab:datastats}. Each encoder sees the same text in the same order
under one tokeniser. Therefore no difference between encoders can originate in
the training text or token segmentation.

The evaluation uses stereotype benchmarks and a capability check,
Table~\ref{tab:datastats}. \emph{CrowS-Pairs} and \emph{StereoSet} measure
association. In both benchmarks, the model is not trained on the benchmark. The
trained encoder is only scored. A consistent preference for the sentence or
completion matching a common stereotype is treated as evidence of lexical
association.

\emph{CrowS-Pairs} contains paired sentences labelled by bias category and by
direction, Table~\ref{tab:datastats}. The dataset label \texttt{race-color}
names its race and colour category. The field \texttt{sent\_more} stores one
member of the sentence pair. The direction label \texttt{stereo} means that
\texttt{sent\_more} is the stereotypical member. The direction label
\texttt{antistereo} means that \texttt{sent\_more} is the anti-stereotypical
member. The direction label affects the official scorer and is analysed in
Section~\ref{sec:scorer}. Shared tokens are tokens that occur in both members of
a sentence pair. An unaligned shared-token pair is a pair for which the
shared-token alignment cannot recover the common token sequence. Category sizes
are uneven, so Section~\ref{sec:category} reports category behaviour rather than
relying only on the aggregate.

\emph{StereoSet} is used in its intrasentence split. Each item supplies a context
sentence with a blank and three candidate fillers: one stereotypical, one
anti-stereotypical, and one unrelated. The unrelated filler tests whether the
encoder prefers a meaningful completion over an unrelated one. The split sizes
are reported in Table~\ref{tab:datastats}.

\emph{WinoBias} is a pronoun-resolution dataset used here as a capability check
rather than a bias measure \citep{zhao2018winobias}. Each item contains
occupation terms and a pronoun that refers to exactly one occupation. Resolving
the pronoun requires sentence understanding. An encoder near chance on this task
cannot support a claim about stereotypical pronoun resolution. The split count
and items per split are in Table~\ref{tab:datastats}.

\begin{table}[t]
\centering
\caption{Concrete examples from the evaluation datasets.}
\label{tab:data}
\footnotesize
\renewcommand{\arraystretch}{1.12}
\begin{tabularx}{\columnwidth}{@{}l>{\raggedright\arraybackslash}X@{}}
\toprule
\textbf{Sentence role} & \textbf{Sample sentence} \\
\midrule
\multicolumn{2}{@{}l@{}}{\textbf{CrowS-Pairs}} \\
stereotypical      & \emph{Cheryl} was bad at saving money and thus had a
                     financial crisis. \\
anti-stereotypical & \emph{John} was bad at saving money and thus had a
                     financial crisis. \\
\midrule
\multicolumn{2}{@{}l@{}}{\textbf{StereoSet}} \\
context            & The chess player was \_\_\_\_. \\
stereotypical      & \ldots \emph{asian}. \\
anti-stereotypical & \ldots \emph{hispanic}. \\
unrelated          & \ldots \emph{fox}. \\
% stereotypical      & \ldots was \emph{asian}. \\
% anti-stereotypical & \ldots was \emph{hispanic}. \\
% unrelated          & \ldots was \emph{fox}. \\
\midrule
\multicolumn{2}{@{}l@{}}{\textbf{WinoBias}} \\
pro-stereotypical  & The developer argued with the designer because
                     \emph{he} did not like the design. \\
anti-stereotypical & \ldots because \emph{she} did not like the design. \\
\bottomrule
\end{tabularx}
\end{table}

\begin{table}[t]
\centering
\caption{Dataset sizes and fixed data choices used in the experiments.}
\label{tab:datastats}
\footnotesize
\begin{adjustbox}{max width=\columnwidth}
\begin{tabular}{llr}
\toprule
Data source & Quantity & Value \\
\midrule
FineWeb-Edu & Training tokens & 7 billion \\
FineWeb-Edu & Sequence length & 128 \\
CrowS-Pairs & Sentence pairs & 1508 \\
CrowS-Pairs & Bias categories & 9 \\
CrowS-Pairs & \texttt{race-color} pairs & 516 \\
CrowS-Pairs & disability pairs & 60 \\
CrowS-Pairs & \texttt{stereo} pairs & 1290 \\
CrowS-Pairs & \texttt{antistereo} pairs & 218 \\
CrowS-Pairs & unaligned shared-token pairs & 8 \\
CrowS-Pairs & mean shared tokens & 15.9 \\
StereoSet & intrasentence items & 2106 \\
StereoSet & bias types & 4 \\
StereoSet & race items & 962 \\
StereoSet & profession items & 810 \\
StereoSet & gender items & 255 \\
StereoSet & religion items & 79 \\
WinoBias & splits & 2 \\
WinoBias & items per split & 374 \\
\bottomrule
\end{tabular}
\end{adjustbox}
\end{table}

\section{Method}
\label{sec:method}

The method links the architecture comparison to item-level measurements. The
computation proceeds from encoder definition, to checkpoint matching, to
benchmark scoring, and then to the capability gate.

\subsection{Architectures}

Architecture is the design choice tested in the study, so all non-architectural
inputs are held fixed where the runs allow. The compared encoders are ordered by
weight reuse in Table~\ref{tab:arch}.

\emph{VanillaBERT} is the unshared control. It follows the BERT encoder design
with independent blocks in sequence \citep{devlin2019bert}. \emph{LoopedBERT}
uses entry blocks, a reused core, and exit blocks. The core blocks are applied
several times in a row. A learned depth marker is added before each pass so that
the reused core can distinguish passes. \emph{ALBERTLoopedBERT} uses one shared
block throughout, following the ALBERT sharing scheme \citep{lan2020albert}.
\emph{HyperloopBERT} keeps the looped layout and carries parallel streams
through the loop. It combines the streams before each core pass, writes the
output back to every stream, and merges the streams before the exit blocks.

All encoders use the same tokeniser, masked language modelling objective,
effective depth, and representation width. Their parameter counts and shared
ratios are shown in Table~\ref{tab:arch}. The shared ratio is the fraction of
effective-depth block applications that reuse block parameters. \emph{HyperloopBERT}
differs from the plain looped design by adding stream mixing and by requiring a
lower peak learning rate for stable training, Table~\ref{tab:reprosettings}.
That encoder is therefore an architecture bundle, meaning a design that changes
more than one mechanism at the same time.

The hyper-connected model carries parallel streams and mixes them at each depth.
The update is
%
\begin{equation}
\mathbf{H}^{(\ell+1)} = \mathbf{A}^{(\ell)}\mathbf{H}^{(\ell)}
                      + \mathbf{B}^{(\ell)} f\!\left(\mathbf{H}^{(\ell)}\right),
\end{equation}
%
where $\mathbf{H}^{(\ell)} \in \mathbb{R}^{n \times d}$ stacks the stream states
at depth $\ell$. Each stream has width $d$. The function $f$ is the shared
transformer block. The matrix $\mathbf{A}^{(\ell)}$ mixes the streams. The
matrix $\mathbf{B}^{(\ell)}$ distributes the block output back across streams.
Both matrices are in $\mathbb{R}^{n \times n}$. The single-stream special case
recovers the ordinary residual connection.

Relative to the hyper-connected looped transformer of
\citet{zeitoun2026hyperloop}, this implementation shares the loop-and-stream
structure and the learned depth-wise mixing. It differs in being trained from
scratch at encoder scale under a masked language modelling objective. The
\emph{HyperloopBERT} results in this paper therefore cover only the
encoder-scale masked-language-model comparison on \emph{CrowS-Pairs} and
\emph{StereoSet} at the matched checkpoints analysed here. They do not cover
every possible hyper-connected looped model.

\begin{table}[t]
\centering
\caption{The four encoders, ordered by how much weight they reuse.}
\label{tab:arch}
\small
\begin{adjustbox}{max width=\columnwidth}
\begin{tabular}{lccc}
\toprule
Architecture & Unique blocks & Parameters & Shared ratio \\
\midrule
VanillaBERT       & 12 & 110.1\,M & 0.00 \\
LoopedBERT        &  6 &  67.6\,M & 0.50 \\
HyperloopBERT     &  6 &  86.5\,M & 0.50 \\
ALBERTLoopedBERT  &  1 &  32.1\,M & 0.92 \\
\bottomrule
\end{tabular}
\end{adjustbox}
\end{table}

\subsection{Matching, scoring, and inference}
\label{sec:rules}

A list of target validation-loss values is fixed before training. The exact
target count and range are in Table~\ref{tab:runfacts}. Validation loss is
measured at regular intervals. A band is a pre-set target-loss level whose
crossing selects a snapshot. The first checkpoint below a target becomes the
snapshot for that band. Encoders are compared only at snapshots assigned to the
same band.

This matching is approximate. Several nominal bands can be crossed inside one
validation interval, so several bands can share a checkpoint. The analysis
therefore uses distinct snapshots, with the per-architecture counts shown in
Table~\ref{tab:runfacts}. Realised losses can also differ slightly inside a
band. Section~\ref{sec:cont} tests whether conclusions remain after continuous
adjustment for realised validation loss.

The sentence score is pseudo-log-likelihood. In each scoring rule, each selected
token position is masked in turn. The probability of the original token is then
evaluated, and the log probabilities are combined. Scoring is performed in the
precision reported in Table~\ref{tab:reprosettings}. Writing $\mathcal{I}$ for
the set of positions a rule scores,
%
\begin{equation}
\mathrm{PLL}(s) = \frac{1}{|\mathcal{I}|}
                  \sum_{i \in \mathcal{I}} \log P\!\left(t_i \mid s_{\setminus i}\right),
\end{equation}
%
where $s_{\setminus i}$ is the sentence with position $i$ masked and $P$ is the
masked language modelling head output. The per-item score is the difference
%
\begin{equation}
e = \mathrm{PLL}(s_{\text{stereo}}) - \mathrm{PLL}(s_{\text{anti}}) .
\end{equation}
%
A positive value means that the stereotypical member receives the higher score.

The \emph{full-sentence} rule scores every position except the special boundary
tokens \texttt{[CLS]} and \texttt{[SEP]}. The token \texttt{[CLS]} marks the
start of the sequence, and \texttt{[SEP]} marks the end or separation boundary.
The \emph{shared-token} rule scores only tokens that the two sentences have in
common. The demographic terms are not scored under the shared-token rule. The
demographic terms remain visible as conditioning context. The benchmark's
released implementation defines the \emph{official} rule. It uses an alignment
that recovers every matching span, sums rather than averages log probabilities,
and assigns the stereotypical role from the direction label. Table~\ref{tab:rules}
summarises these rules.

One \emph{CrowS-Pairs} item illustrates the difference. The pair
``\ldots{}and whether he would come forward'' against ``\ldots{}and whether she
would come forward'' is labelled \texttt{antistereo}, Table~\ref{tab:datastats}.
In that item, \texttt{sent\_more} is the anti-stereotypical sentence. The
full-sentence and shared-token rules treat \texttt{sent\_more} as stereotypical.
The official rule reverses the role for that item because of the direction
label. The same pair can therefore contribute with the opposite sign under the
official rule.

The contrast is paired by benchmark item. Every item is scored under the
baseline and another encoder at the same matched point. The paired difference is
$d_i = e_i^{\text{base}} - e_i^{\text{other}}$. The contrast $\Delta$ is the
mean of those differences over the benchmark items. A positive contrast means
that \emph{VanillaBERT} records the larger stereotype effect.

The null hypothesis $\mathbb{E}[d_i] = 0$ is tested by a paired sign-flip
permutation test. The draw count and bootstrap count are in
Table~\ref{tab:reprosettings}. The $p$-value uses the Phipson--Smyth estimate
\citep{phipson2010permutation}. That estimate avoids a zero permutation
$p$-value when a finite number of random sign flips is sampled. Confidence
intervals are percentile item bootstraps, so the intervals describe item
sampling only. Where several contrasts are formed against the same baseline,
Holm's step-down procedure controls the family-wise error rate within that
family.

\begin{table}[t]
\centering
\caption{The three scoring rules compared in this paper.}
\label{tab:rules}
\footnotesize
\begin{adjustbox}{max width=\columnwidth}
\begin{tabular}{lccc}
\toprule
Rule & Tokens scored & Combined by & Stereotypical member \\
\midrule
Full-sentence  & all             & average & \texttt{sent\_more} always \\
Shared-token   & words in common & average & \texttt{sent\_more} always \\
Official       & words in common & sum     & by direction label \\
\bottomrule
\end{tabular}
\end{adjustbox}
\end{table}

\label{sec:gate}
The capability gate defines when a stereotype measurement can be interpreted as
more than a lexical association score. GLUE is a general language understanding
benchmark used in the planned gate \citep{wang2018glue}. The pre-specified gate
required above-chance performance on GLUE, a baseline stereotype-effect interval
excluding zero, and above-chance accuracy on the pro-stereotype split of
\emph{WinoBias}. The \emph{WinoBias} component tests whether the encoder
resolves the pronoun at all.

For \emph{WinoBias}, the pronoun token is replaced with \texttt{[MASK]}. The
masked language modelling head scores the candidate pronouns from the item. A
prediction is correct when the higher-scoring pronoun matches the gold referent
encoded by the dataset item. Accuracy is the share of correct items in each
split.

\section{Experimental Setup}
\label{sec:setup}

The experimental setup fixes training data and most hyper-parameters across
architectures. It also records the exceptions needed to make the runs usable.
Table~\ref{tab:runfacts} gives the run-level facts that affect interpretation,
and Table~\ref{tab:reprosettings} gives implementation settings.

The encoders reach the lowest-loss common matched point after different amounts
of training, Table~\ref{tab:runfacts}. Matching on validation loss therefore
does not match the training budget. \emph{HyperloopBERT} is the only architecture
with a different peak learning rate, Table~\ref{tab:reprosettings}. At the
common rate its loss became non-finite. Holding the rate fixed would have
compared stable models with a failed run. The reported \emph{HyperloopBERT}
comparisons therefore treat that model as a stability-adjusted architecture
bundle. They are not an isolated test of streams. Each encoder was trained once,
Table~\ref{tab:reprosettings}. Reported intervals quantify benchmark-item
uncertainty, not variation across training runs.

\begin{table}[t]
\centering
\caption{Run facts used to interpret matched checkpoints and training stability.}
\label{tab:runfacts}
\footnotesize
\begin{adjustbox}{max width=\columnwidth}
\begin{tabular}{llr}
\toprule
Quantity & Scope & Value \\
\midrule
Target validation-loss bands & pre-set list & 7 \\
Highest target loss & pre-set list & 5.0 \\
Lowest target loss & pre-set list & 2.2 \\
Distinct snapshots & all encoders & 21 \\
Distinct snapshots & VanillaBERT & 4 \\
Distinct snapshots & LoopedBERT & 5 \\
Distinct snapshots & ALBERTLoopedBERT & 5 \\
Distinct snapshots & HyperloopBERT & 7 \\
Training tokens to deepest point & minimum & 1.53 billion \\
Training tokens to deepest point & maximum & 2.07 billion \\
Initial flat-loss region exit & HyperloopBERT lower-rate run & 0.5 billion \\
Initial flat-loss region exit & VanillaBERT and LoopedBERT runs & 1.4 billion \\
HyperloopBERT unstable interval start loss & common-rate run & 2.53 \\
HyperloopBERT unstable interval end loss & common-rate run & 6.86 \\
ALBERTLoopedBERT spike recovery & optimiser steps & about 1000 \\
\bottomrule
\end{tabular}
\end{adjustbox}
\end{table}

\section{Results}
\label{sec:results}

The results first establish the single-point comparison and then test whether
the comparison changes under specified analysis choices. The final paragraphs
state the capability-gate outcome, because that outcome limits every bias
interpretation in the section.

\subsection{Single-point comparison and concrete predictions}

The single-point analysis uses the lowest-loss matched point common to all
encoders. Figure~\ref{fig:training} shows why matching is needed before that
point is interpreted. The validation-loss curves for \emph{VanillaBERT} and
\emph{LoopedBERT} nearly overlap, Figure~\ref{fig:training}.
\emph{ALBERTLoopedBERT} remains higher, which indicates lower
language-modelling quality in these runs. \emph{HyperloopBERT} leaves the early
flat-loss region earlier but later diverges to non-finite loss. The instability
values and snapshot counts are reported in Table~\ref{tab:runfacts}. Every
snapshot analysed below predates the non-finite values.

\begin{figure}[t]
\centering
\includegraphics[width=0.85\columnwidth]{images/fig2_training_dynamics.pdf}
\Description{Validation loss plotted against training tokens for four encoders.
HyperloopBERT descends earliest and its curve ends in a divergence marker near
two billion tokens. ALBERTLoopedBERT shows one spike that returns to trend.
VanillaBERT and LoopedBERT overlap closely and end lowest.}
\caption{Validation loss against training tokens for the four encoders.}
\label{fig:training}
\end{figure}

Table~\ref{tab:band} reports realised validation loss, mean item effect under
the full-sentence rule, percentile item-bootstrap interval, contrast against
\emph{VanillaBERT}, and Holm-corrected permutation $p$-value. In
Table~\ref{tab:band}, \emph{LoopedBERT} and the \emph{HyperloopBERT} bundle have
lower \emph{CrowS-Pairs} effects than the unshared baseline. Both contrasts
survive Holm correction. \emph{ALBERTLoopedBERT} does not show a corrected
contrast against the baseline.

The two looped variants share the same loop and differ by the stream mechanism
and training-rate adjustment. Their contrasts against \emph{VanillaBERT} are
similar at the lowest-loss matched point, Table~\ref{tab:band}. No direct
\emph{HyperloopBERT}-versus-\emph{LoopedBERT} test is reported.

\begin{table}[t]
\centering
\caption{\emph{CrowS-Pairs} results at the lowest-loss matched point.}
\label{tab:band}
\footnotesize
\begin{adjustbox}{max width=\columnwidth}
\begin{tabular}{lccccc}
\toprule
Architecture & Loss & Effect & 95\% interval & $\Delta$ & $p_{\text{Holm}}$ \\
\midrule
VanillaBERT      & 2.183 & 0.0707 & [0.045, 0.097] & --- & --- \\
LoopedBERT       & 2.192 & 0.0471 & [0.022, 0.072] & 0.0237 & 0.0003 \\
ALBERTLoopedBERT & 2.163 & 0.0591 & [0.034, 0.085] & 0.0116 & 0.0653 \\
HyperloopBERT    & 2.196 & 0.0466 & [0.022, 0.071] & 0.0241 & 0.0003 \\
\bottomrule
\end{tabular}
\end{adjustbox}
\end{table}

\label{sec:qualitative}
The aggregate in Table~\ref{tab:band} can be interpreted through direct model
outputs. Table~\ref{tab:qual} reports hand-built masked-pronoun probes scored at
the same matched snapshots. These probes are not part of any benchmark. The
first probe is ``The scientist said that \texttt{[MASK]} was ready for the
lab.'' The other probes replace ``scientist'' with ``teacher'' and
``secretary.'' On the scientist sentence, the unshared baseline assigns higher
probability to \emph{he} than to \emph{she}, Table~\ref{tab:qual}. On the
teacher sentence, the same encoder prefers \emph{she}. The direction is
therefore occupation-specific, not a fixed preference for one pronoun.

The secretary probe shows why paired differences are used. The common stereotype
would favour \emph{she}, but every encoder assigns higher probability to
\emph{he}, Table~\ref{tab:qual}. Pronoun frequency can therefore compete with an
occupational association. The paired construction reduces this general
pronoun-frequency component by changing only the group term inside a sentence.

\begin{table}[t]
\centering
\caption{Probability assigned to \emph{he} and \emph{she} at the masked pronoun
of three probes.}
\label{tab:qual}
\footnotesize
\begin{adjustbox}{max width=\columnwidth}
\begin{tabular}{lccc}
\toprule
 & \emph{scientist} & \emph{teacher} & \emph{secretary} \\
Encoder & $P(\textit{he})/P(\textit{she})$ & & \\
\midrule
VanillaBERT      & \textbf{0.305 / 0.005} & \textbf{0.009 / 0.052} & 0.115 / 0.006 \\
LoopedBERT       & 0.410 / 0.026 & 0.001 / 0.001 & 0.308 / 0.023 \\
ALBERTLoopedBERT & 0.637 / 0.017 & 0.007 / 0.005 & 0.155 / 0.004 \\
HyperloopBERT    & 0.610 / 0.033 & 0.001 / 0.001 & 0.066 / 0.004 \\
\bottomrule
\end{tabular}
\end{adjustbox}
\end{table}

\subsection{Sensitivity analyses}
\label{sec:allband}

The single-point result motivates checks of the analysis choices. Each check
asks whether the lower \emph{CrowS-Pairs} effect for \emph{LoopedBERT} and the
\emph{HyperloopBERT} bundle remains when one choice changes.

Changing the matched point changes the comparison. Figure~\ref{fig:delta} plots
the \emph{CrowS-Pairs} contrast against realised validation loss. The
baseline-versus-looped contrast changes sign across matched points. The
reduction in Table~\ref{tab:band} appears only at the lowest-loss looped point.
The other looped points have the opposite sign. The contrasts against the other
weight-reusing encoders also vary by matched point, Figure~\ref{fig:delta}. The
\emph{CrowS-Pairs} matched-point contrasts against \emph{VanillaBERT} are
exploratory and repeated. The matched points within an encoder are successive
checkpoints of one run, so the points are not independent.

\begin{figure}[t]
\centering
\includegraphics[width=0.85\columnwidth]{images/fig3_delta_across_loss.pdf}
\Description{Architecture contrast plotted against matched validation loss with
bootstrap error bars. All three contrast curves cross the zero line, and the
baseline-versus-looped curve is negative at four of its five points and positive
only at the lowest-loss point.}
\caption{\emph{CrowS-Pairs} contrast against the unshared baseline at every
matched point.}
\label{fig:delta}
\end{figure}

\label{sec:cont}
The matched-point variation motivates a continuous adjustment for realised
capability. The regression uses \emph{CrowS-Pairs} item effects as the response,
architecture as a categorical factor with \emph{VanillaBERT} as reference, and
mean-centred realised loss as a covariate. The model is fitted by least squares.
Architecture and loss vary only across snapshots, so the identifying units are
snapshots rather than item rows. A wild cluster bootstrap is an inference check
that resamples snapshot-level residual patterns while preserving dependence
among item rows from the same snapshot. Table~\ref{tab:model} reports the
snapshot-level fit and that bootstrap.

After adjustment for realised validation loss, no architecture coefficient in
the \emph{CrowS-Pairs} item-effect regression is distinguishable from zero at
the significance rule used for bold entries, Table~\ref{tab:model}. Realised
loss remains associated with the mean \emph{CrowS-Pairs} effect,
Table~\ref{tab:model}. The loss coefficient is on the same scale as the
contrasts in Table~\ref{tab:band}. Remaining capability differences are
therefore large enough to affect the single-point reading. This is not evidence
that architecture has no effect. With the snapshot count in Table~\ref{tab:runfacts}
and the seed count in Table~\ref{tab:reprosettings}, the analysis has limited
power for small architecture effects.

\begin{table}[t]
\centering
\caption{\emph{CrowS-Pairs} effect explained by encoder design and by training
progress, over the 21 snapshots.}
\label{tab:model}
\footnotesize
\begin{adjustbox}{max width=\columnwidth}
\begin{tabular}{lccc}
\toprule
Term & Coefficient & $p$ (snapshot) & $p$ (bootstrap) \\
\midrule
LoopedBERT        & $+0.0047$ & 0.424 & 0.488 \\
ALBERTLoopedBERT  & $-0.0065$ & 0.280 & 0.169 \\
HyperloopBERT     & $-0.0066$ & 0.265 & 0.191 \\
Realised loss     & $-0.0184$ & \textbf{0.0002} & \textbf{0.0049} \\
\bottomrule
\end{tabular}
\end{adjustbox}
\end{table}

\label{sec:scorer}
The loss-adjusted result leaves the scorer as a separate source of variation.
The shared-token and official \emph{CrowS-Pairs} rules rescore the same items at
the lowest-loss matched snapshots. Table~\ref{tab:rules2} gives the exact
contrasts and the count of contrasts that survive Holm correction. Neither
alternative rule preserves the corrected result. Under the shared-token rule, no
contrast survives Holm correction. Under the official rule, no contrast survives
correction, and the \emph{LoopedBERT} contrast reverses sign,
Table~\ref{tab:rules2}.

The official rule differs from the shared-token rule in token alignment,
aggregation, and direction convention. The decomposition in
Table~\ref{tab:rules2} fixes the official token set and separates aggregation
from direction. Changing from averaging to summing changes scale. Changing the
direction convention changes sign for the \emph{LoopedBERT} and
\emph{ALBERTLoopedBERT} contrasts. The sign reversal is therefore attributable
to the benchmark direction convention for the \texttt{antistereo} subset, whose
size is given in Table~\ref{tab:datastats}.

The released BERT checkpoint \texttt{bert-base-uncased} was rescored with this
implementation. The absolute score differs from the value reported by
\citet{nangia2020crows}. The public checkpoint and current scoring library are
not the same execution environment as the original benchmark release. All four
encoders in this paper are scored with the same code and the same environment,
so the architecture contrasts use a common scorer. The shared-token
prefix-and-suffix alignment fails on the unaligned subset shown in
Table~\ref{tab:datastats}. These pairs are instances of the item-quality problem
documented by \citet{blodgett2021salmon}.

\begin{table}[t]
\centering
\caption{Contrast against \emph{VanillaBERT} on \emph{CrowS-Pairs} at the
lowest-loss matched point, under each scoring rule.}
\label{tab:rules2}
\footnotesize
\begin{adjustbox}{max width=\columnwidth}
\begin{tabular}{lccccc}
\toprule
Scoring rule & $N$ & vs LoopedBERT & vs ALBERTLoopedBERT & vs HyperloopBERT & Sig. \\
\midrule
Full-sentence (mean, all)      & 1508 & \textbf{0.0237} & 0.0116    & \textbf{0.0241} & 2 of 3 \\
Shared-token (mean, shared)    & 1500 & 0.0054          & 0.0062    & 0.0143          & 0 of 3 \\
Official (sum, shared, signed) & 1508 & $-0.0519$       & $-0.0163$ & $+0.0488$       & 0 of 3 \\
\addlinespace[2pt]
\multicolumn{6}{@{}l}{\emph{official token set and sum, direction convention removed}} \\
\texttt{sent\_more} stereotypical & 1508 & $+0.0919$ & $+0.0948$ & $+0.1305$ & 0 of 3 \\
\addlinespace[2pt]
\multicolumn{6}{@{}l}{Sig. counts Holm-surviving contrasts; bold marks entries passing the stated decision rule.} \\
\bottomrule
\end{tabular}
\end{adjustbox}
\end{table}

\label{sec:stereoset}
A second benchmark tests whether the single-point direction is specific to
\emph{CrowS-Pairs}. The same lowest-loss matched snapshots are scored on the
\emph{StereoSet} intrasentence split. The \emph{StereoSet} paired effect is the
pseudo-log-likelihood difference between the stereotypical and
anti-stereotypical fillers in the blank position. Table~\ref{tab:stereoset}
reports the paired effect, the benchmark stereotype score \emph{SS}, the
language-modelling score \emph{LMS}, the contrast, and the Holm-corrected
$p$-value. \emph{SS} is the benchmark percentage favouring the stereotypical
completion. \emph{LMS} is the benchmark percentage favouring a meaningful
completion over an unrelated one.

The direction agrees with the \emph{CrowS-Pairs} single-point result for all
weight-reusing encoders, Table~\ref{tab:stereoset}. The inferential conclusion
does not agree. On the \emph{StereoSet} intrasentence items, only the paired
effect for \emph{HyperloopBERT} versus \emph{VanillaBERT} survives Holm
correction at the lowest-loss matched point, Table~\ref{tab:stereoset}. On
\emph{CrowS-Pairs}, the corrected contrasts were against \emph{LoopedBERT} and
\emph{HyperloopBERT}, Table~\ref{tab:band}. \emph{CrowS-Pairs} and
\emph{StereoSet} therefore flag different encoders after correction.

\begin{table}[t]
\centering
\caption{Results on the second benchmark, \emph{StereoSet}, at the same matched
point.}
\label{tab:stereoset}
\footnotesize
\begin{adjustbox}{max width=\columnwidth}
\begin{tabular}{lccccc}
\toprule
Architecture & Effect & SS & LMS & $\Delta$ & $p_{\text{Holm}}$ \\
\midrule
VanillaBERT      & 0.0458 & 51.4 & 86.0 & --- & --- \\
LoopedBERT       & 0.0310 & 50.6 & 86.3 & 0.0148 & 0.084 \\
ALBERTLoopedBERT & 0.0294 & 50.9 & 87.1 & 0.0164 & 0.084 \\
HyperloopBERT    & 0.0158 & 49.7 & 86.3 & 0.0300 & \textbf{0.0009} \\
\bottomrule
\end{tabular}
\end{adjustbox}
\end{table}

\label{sec:category}
The benchmark difference motivates checking whether the aggregate hides category
variation. Figure~\ref{fig:robust} shows that category contrasts appear on both
sides of zero. Several categories reverse sign depending on which
weight-reusing encoder is compared with the baseline. The largest category and
smallest category sizes are reported in Table~\ref{tab:datastats}. These
analyses are descriptive. The category analyses followed the aggregate
instability and are not confirmatory tests of a pre-stated category hypothesis.

\begin{figure*}[t]
\centering
\includegraphics[width=0.48\textwidth,trim=330 0 0 0,clip]{images/fig4_scorer_and_category.pdf}
\Description{Category contrasts for nine bias categories, with values on both
sides of zero.}
\caption{\emph{CrowS-Pairs} contrasts by bias category.}
\label{fig:robust}
\end{figure*}

\label{sec:capability}
The capability gate limits the interpretation of all preceding effects. On
\emph{WinoBias}, accuracies for every encoder and split lie just above and just
below chance, Table~\ref{tab:gate}. An encoder that cannot resolve the pronoun
cannot support a claim about stereotypical preference in pronoun resolution. The
GLUE component was not run because the failed coreference component already made
the gate fail. The baseline-interval component is met by the \emph{VanillaBERT}
row of Table~\ref{tab:band}.

The failed \emph{WinoBias} check limits the reported effects to lexical
association claims. The reported effects measure lexical association in the
masked language modelling head. The reported effects do not show that a ranker
or classifier fine-tuned from \emph{VanillaBERT}, \emph{LoopedBERT},
\emph{ALBERTLoopedBERT}, or \emph{HyperloopBERT} would treat demographic groups
differently in use.

\begin{table}[t]
\centering
\caption{Capability-gate outcomes at the evaluated checkpoints.}
\label{tab:gate}
\footnotesize
\begin{adjustbox}{max width=\columnwidth}
\begin{tabular}{llr}
\toprule
Gate component & Quantity & Value \\
\midrule
WinoBias & minimum accuracy & 0.489 \\
WinoBias & maximum accuracy & 0.503 \\
WinoBias & chance level & 0.5 \\
WinoBias & items per split & 374 \\
GLUE & run status & not run \\
CrowS-Pairs baseline interval & gate status & met \\
Overall capability gate & gate status & failed \\
\bottomrule
\end{tabular}
\end{adjustbox}
\end{table}

\section{Discussion}
\label{sec:discussion}

The results separate a single matched-point finding from a general architecture
claim. At the lowest-loss \emph{CrowS-Pairs} point, \emph{LoopedBERT} and the
\emph{HyperloopBERT} bundle are lower than \emph{VanillaBERT},
Table~\ref{tab:band}. That corrected contrast does not repeat across matched
points or alternative scorers, Figure~\ref{fig:delta} and Table~\ref{tab:rules2}.
The pipeline checks used the same data order, scorer code, and snapshot-selection
rule for all encoders. The remaining issue is dependence on the matched point
and scorer.

The matched-point dependence suggests that residual capability still matters.
After continuous adjustment for realised validation loss, no architecture term
is distinguishable from zero at the stated decision rule, Table~\ref{tab:model}.
Validation loss remains associated with the \emph{CrowS-Pairs} item effect,
Table~\ref{tab:model}. This association is not a causal proof. The snapshot
count is limited, and all checkpoints come from the seed count reported in
Table~\ref{tab:reprosettings}.

The scorer dependence has a different source. The three scoring rules make
different choices about token sets, aggregation, and direction labels,
Table~\ref{tab:rules}. The decomposition shows that the sign reversal follows
from the direction convention for \texttt{antistereo} pairs,
Table~\ref{tab:rules2}. Treating \texttt{sent\_more} as stereotypical and using
the dataset direction label are both explicit conventions. A conclusion that
holds under only one convention is therefore not enough in this study to
attribute the measured difference to architecture alone.

Retrieval and classification pipelines often choose among pre-trained encoders.
In this study, an architecture-level association comparison changes under
matched capability level, scoring rule, and direction convention. The comparison
therefore supports model-selection claims only when the matched point, scoring
rule, and direction convention are stated with the result.

\section{Limitations}
\label{sec:limits}

Several limits affect how the results should be read. The most important limit
is that training uses the seed count shown in Table~\ref{tab:reprosettings}. The
runs are therefore not replications. Every interval reported here bootstraps
over benchmark items and quantifies item sampling alone. The intervals do not
quantify variation between training runs. The ordering of encoders may change
under re-training with another seed. The single-seed design also means that the
loss trajectories in Figure~\ref{fig:training} cannot separate architecture from
run-specific optimisation variation.

The analyses across matched points, continuous loss adjustment, scoring-rule
decomposition, and category comparisons were performed after the single-point
result had been examined. The analyses are exploratory. Only the lowest-loss
point contrasts against \emph{VanillaBERT} were specified in advance. The
reported corrected results should therefore be read as measurements from this
run, not as confirmatory evidence for a general architecture effect.

The encoders have not converged at the lowest-loss matched point. The validation
losses in Table~\ref{tab:band} are still changing, and Figure~\ref{fig:training}
does not show a finished training curve. Whether the architecture contrasts
would differ for fully trained encoders is untested here. \emph{HyperloopBERT}
also needed a lower peak learning rate, Table~\ref{tab:reprosettings}. Its
comparison therefore bundles the stream architecture with a stability adjustment.

Only two stereotype benchmarks are used, and both are correlation-type
instruments in the sense of \citet{wang2025difference}. Both are scored from
the same masked language modelling head. Agreement between these instruments is
therefore weaker evidence than agreement between independent measurement
approaches would be. The capability gate fails, Table~\ref{tab:gate}, so no
statement about downstream behaviour is supported. Both instruments also have
documented item-quality problems \citep{blodgett2021salmon}. The unaligned
\emph{CrowS-Pairs} subset in Table~\ref{tab:datastats} is one instance.

The scoring-rule decomposition isolates the direction convention at one matched
point. The decomposition does not prove that the same attribution holds at other
matched points or on \emph{StereoSet}. Those settings were not rescored under
the official rule. The category analysis is also descriptive. Category sizes are
uneven, Table~\ref{tab:datastats}, and no category hypothesis was specified
before examining the aggregate.

The evaluation is English-only. An India-centric instrument covering caste and
religion \citep{khandelwal2023indianbhed} was planned, but the available mirror
could not be used with the English-only tokeniser. Extending the protocol to
that setting is the first follow-up direction.

\section{Reproducibility}
\label{sec:repro}

A reader who reruns the study needs the training corpus order, model
definitions, saved checkpoints, scorer code, and inference settings. The fixed
data choices are in Section~\ref{sec:data}. The encoder definitions are in
Table~\ref{tab:arch}. The training and analysis settings are listed in
Table~\ref{tab:reprosettings}.

All runs verify a checksum and block-count manifest before training. Attention
uses the memory-efficient exact implementation of \citet{dao2022flashattention},
and the run log confirms that implementation was active throughout. Pre-training
uses AdamW, bfloat16 autocast, and H100 GPUs, with exact settings in
Table~\ref{tab:reprosettings}. Snapshots are selected by first crossing of
pre-set validation-loss targets, Table~\ref{tab:runfacts}. The \emph{StereoSet}
replication and official-rule rescoring both rerun on CPU from the same
snapshots. Re-running the released scripts reproduces the stored scores within
the CPU rerun tolerance in Table~\ref{tab:reprosettings}.
% CAMERA-READY: replace the anonymised mirror with the permanent repository URL.
Code, per-item scores for every snapshot, the masked-position predictions behind
Table~\ref{tab:qual}, and the analysis scripts are available at
\url{https://anonymous.4open.science/r/HyperloopBert/}.

\begin{table}[t]
\centering
\caption{Training, scoring, and analysis settings needed to reproduce the study.}
\label{tab:reprosettings}
\footnotesize
\begin{adjustbox}{max width=\columnwidth}
\begin{tabular}{llr}
\toprule
Setting & Scope & Value \\
\midrule
Peak learning rate & main encoders & $3 \times 10^{-4}$ \\
Peak learning rate & HyperloopBERT & $1.5 \times 10^{-4}$ \\
Batch size & pre-training & 512 \\
Warmup share & pre-training & 10\% \\
Gradient clipping & pre-training & 1.0 \\
Autocast format & pre-training & bfloat16 \\
Random seeds & pre-training & 1 \\
Seed value & pre-training & 42 \\
GPU type & pre-training & H100 \\
Scoring precision & benchmark scoring & 32-bit \\
Permutation draws & sign-flip test & $10^4$ \\
Bootstrap resamples & item intervals & $10^4$ \\
Wild bootstrap draws & cluster bootstrap & 9999 \\
Python & analysis & 3.10 \\
NumPy & analysis & 2.2 \\
pandas & analysis & 2.3 \\
SciPy & analysis & 1.15 \\
statsmodels & analysis & 0.14 \\
CPU rerun tolerance & stored scores & $10^{-6}$ \\
\bottomrule
\end{tabular}
\end{adjustbox}
\end{table}

\section{Conclusion}
\label{sec:conclusion}

The study began from a practical question about encoder choice. If one
weight-reusing encoder appears to have lower stereotype association than an
unshared encoder, the comparison should remain similar under nearby matched
training points and stated scoring rules.

This study compared encoders that differ in weight reuse across depth, while
holding data, tokenisation, objective, representation width, and effective depth
fixed. At the lowest-loss matched point on \emph{CrowS-Pairs}, \emph{LoopedBERT}
and the \emph{HyperloopBERT} bundle show lower effects than \emph{VanillaBERT}.
Both contrasts survive Holm correction, Table~\ref{tab:band}. That result does
not hold across the subsequent checks. The looped result appears at only one
matched point, Figure~\ref{fig:delta}. The architecture terms become
indistinguishable from zero after continuous loss adjustment,
Table~\ref{tab:model}. The result survives correction under neither alternative
scoring rule, Table~\ref{tab:rules2}. The official rule reverses the
\emph{LoopedBERT} and \emph{ALBERTLoopedBERT} contrasts because of the direction
convention for \texttt{antistereo} pairs.

The second benchmark gives a related but not identical result. On
\emph{StereoSet}, the direction agrees for all weight-reusing encoders, but only
the \emph{HyperloopBERT} contrast remains corrected, Table~\ref{tab:stereoset}.
The category analysis and capability gate further limit the interpretation,
Figure~\ref{fig:robust} and Table~\ref{tab:gate}. The evidence supports a
measurement recommendation. Bias comparisons should report matched capability
levels, scoring rules, and direction conventions before attributing a measured
difference to encoder design.

\section*{Declaration on Generative AI}

The writing process included a final language-editing step. ChatGPT was used to
enhance grammar, rectify spelling mistakes, and rephrase writing. All content
was reviewed and verified by the authors, who take full responsibility for it.

\bibliographystyle{ACM-Reference-Format}
\bibliography{reference}

\end{document}